\begin{abstract}
La fusión multimodal de datos hiperespectrales (HSI) y de radar de apertura sintética (SAR) representa un avance clave en percepción remota para monitoreo crítico de desastres, agricultura y vigilancia ambiental. Este artículo presenta una arquitectura basada en redes neuronales profundas que integra características espectrales ricas de HSI con información estructural y de penetración de SAR mediante mecanismos de atención cruzada y fusión asimétrica. Se revisa el estado del arte en fusión multimodal, se detalla la propuesta metodológica con módulos de extracción multi-escala y transformers, y se validan resultados en conjuntos de datos públicos como Augsburg. Los experimentos demuestran mejoras significativas en precisión de clasificación y robustez bajo condiciones adversas (nubes, noche). Se discuten limitaciones computacionales y se proponen direcciones futuras hacia modelos ligeros y en tiempo real. (148 palabras)
\end{abstract}

\begin{IEEEkeywords}
Cross-Attention Mechanism; Deep Neural Networks; Feature Extraction; High-Fidelity Remote Sensing; Hyperspectral Imaging (HSI); Multimodal Data Fusion; Remote Sensing Classification; Synthetic Aperture Radar (SAR); Transformers.
\end{IEEEkeywords}